% ═══════════════════════════════════════════════════════════════
% SPANISCH ARBEITSBLATT - DS-01: Mi nueva casa
% Klasse 9, Unidad 2
% ═══════════════════════════════════════════════════════════════

\documentclass[10pt, a4paper]{article}

\usepackage[utf8]{inputenc}
\usepackage[T1]{fontenc}
\usepackage[ngerman]{babel}

% --- SCHRIFTEN ---
\usepackage{palatino}
\usepackage{helvet}
\newcommand{\sanstext}[1]{{\fontfamily{phv}\selectfont #1}}

\usepackage{microtype}
\usepackage[margin=1.4cm, top=1.6cm, bottom=1.3cm]{geometry}
\usepackage{enumitem}
\usepackage{tabularx}
\usepackage{booktabs}
\usepackage{array}
\usepackage{multicol}
\usepackage{tikz}
\usetikzlibrary{calc}
\usepackage{amssymb}

\usepackage{fancyhdr}
\usepackage{lastpage}
\usepackage{needspace}

\usepackage{xcolor}
\usepackage{tcolorbox}
\tcbuselibrary{skins}

% ═══════════════════════════════════════════════════════════════
% FARBEN
% ═══════════════════════════════════════════════════════════════

\definecolor{kopfzeile}{HTML}{1A1A1A}
\definecolor{akzent}{HTML}{C62828}
\definecolor{hellgrau}{HTML}{F2F2F2}
\definecolor{mittelgrau}{HTML}{777777}
\definecolor{dunkelgrau}{HTML}{444444}

% ═══════════════════════════════════════════════════════════════
% VARIABLEN
% ═══════════════════════════════════════════════════════════════

\newcommand{\thema}{Mi nueva casa}
\newcommand{\klasse}{9}
\newcommand{\maxpunkte}{30}

% ═══════════════════════════════════════════════════════════════
% KOPF- UND FUSSZEILE
% ═══════════════════════════════════════════════════════════════

\pagestyle{fancy}
\fancyhf{}
\renewcommand{\headrulewidth}{0pt}
\renewcommand{\footrulewidth}{0pt}
\cfoot{\sanstext{\footnotesize\textcolor{mittelgrau}{\thepage\,/\,\pageref{LastPage}}}}

% ═══════════════════════════════════════════════════════════════
% BOXEN
% ═══════════════════════════════════════════════════════════════

\newtcolorbox{grammatikbox}{
    colback=hellgrau,
    colframe=hellgrau,
    boxrule=0pt,
    arc=0pt,
    left=8pt, right=8pt, top=8pt, bottom=6pt,
    borderline north={1.5pt}{0pt}{dunkelgrau}
}

\newtcolorbox{vokabelbox}{
    colback=white,
    colframe=mittelgrau,
    boxrule=0.5pt,
    arc=0pt,
    left=6pt, right=6pt, top=6pt, bottom=6pt
}

\newcommand{\regel}[1]{%
    \tikz[baseline=(X.base)]\node[fill=akzent, text=white,
    font=\sanstext\scriptsize\bfseries, inner sep=2pt, rounded corners=1pt](X){#1};%
}

% ═══════════════════════════════════════════════════════════════
% BEFEHLE
% ═══════════════════════════════════════════════════════════════

\newcommand{\es}[1]{\textbf{#1}}
\newcommand{\de}[1]{{\small\textit{\textcolor{mittelgrau}{#1}}}}
\newcommand{\luecke}[1]{\rule{#1}{0.4pt}}
\newcommand{\punktebox}[1]{\hfill\sanstext{\small[\_\_/#1]}}

\newcommand{\aufgabe}[2]{%
    \needspace{4\baselineskip}%
    \vspace{10pt}%
    \noindent\sanstext{\textbf{#1}\quad#2}%
    \vspace{5pt}%
    \nopagebreak[4]%
}

\newlist{teil}{enumerate}{1}
\setlist[teil]{label=\alph*), leftmargin=1.8em, itemsep=3pt, parsep=0pt, topsep=3pt}

\setlength{\parindent}{0pt}
\setlength{\parskip}{4pt}
\setlength{\columnsep}{20pt}

% ═══════════════════════════════════════════════════════════════
% DOKUMENT
% ═══════════════════════════════════════════════════════════════

\begin{document}

% ═══════════════════════════════════════════════════════════════
% HEADER
% ═══════════════════════════════════════════════════════════════

\noindent
\begin{tikzpicture}[remember picture, overlay]
    \fill[kopfzeile] (current page.north west) rectangle +(\paperwidth, -1cm);
    \node[anchor=west, text=white, font=\sanstext\large\bfseries]
        at ([xshift=1.4cm, yshift=-0.5cm]current page.north west) {\thema};
    \node[anchor=east, text=white, font=\sanstext\small]
        at ([xshift=-1.4cm, yshift=-0.5cm]current page.north east)
        {Spanisch\,·\,Klasse \klasse\,·\,Arbeitsblatt};
\end{tikzpicture}

\vspace{0.5cm}

\noindent
\sanstext{\small\textbf{Name:}} \luecke{5cm} \hfill
\sanstext{\small\textbf{Datum:}} \luecke{2.5cm}

\vspace{6pt}

% ═══════════════════════════════════════════════════════════════
% KASTEN V1: VOKABELN
% ═══════════════════════════════════════════════════════════════

\begin{vokabelbox}
\sanstext{\footnotesize\textbf{Kasten V1: Vocabulario}}

\vspace{4pt}

\begin{multicols}{2}

{\small\textbf{Partes de la casa} \de{Hausräume}}

\vspace{2pt}

{\footnotesize
\begin{tabular}{@{}ll@{}}
\es{el pasillo} & \de{der Flur} \\
\es{el dormitorio} & \de{das Schlafzimmer} \\
\es{el salón} & \de{das Wohnzimmer} \\
\es{la cocina} & \de{die Küche} \\
\es{el balcón} & \de{der Balkon} \\
\es{el baño} & \de{das Badezimmer} \\
\es{el aseo} & \de{die Toilette} \\
\end{tabular}}

\columnbreak

{\small\textbf{Muebles y objetos} \de{Möbel und Gegenstände}}

\vspace{2pt}

{\footnotesize
\begin{tabular}{@{}ll@{}}
\es{la lámpara} & \de{die Lampe} \\
\es{la cama} & \de{das Bett} \\
\es{el armario} & \de{der Schrank} \\
\es{la estantería} & \de{das Regal} \\
\es{la planta} & \de{die Pflanze} \\
\es{el sillón} & \de{der Sessel} \\
\es{el sofá} & \de{das Sofa} \\
\es{el cuadro} & \de{das Bild} \\
\es{el televisor} & \de{der Fernseher} \\
\es{la nevera} & \de{der Kühlschrank} \\
\es{el cubo de basura} & \de{der Mülleimer} \\
\es{el reloj} & \de{die Uhr} \\
\es{el horno} & \de{der Ofen} \\
\es{la ducha} & \de{die Dusche} \\
\es{el espejo} & \de{der Spiegel} \\
\es{la guitarra} & \de{die Gitarre} \\
\es{el póster} & \de{das Poster} \\
\end{tabular}}

\end{multicols}
\end{vokabelbox}

\vspace{4pt}

% ═══════════════════════════════════════════════════════════════
% KASTEN G1: BEWEGUNGSVERBEN
% ═══════════════════════════════════════════════════════════════

\begin{grammatikbox}
\sanstext{\footnotesize\textbf{Kasten G1: Los verbos de movimiento}}

\vspace{4pt}

\begin{multicols}{2}

{\small Ergänze die Präpositionen:}

\vspace{3pt}

{\footnotesize
\begin{tabular}{@{}cll@{}}
$\downarrow$ & \es{bajar} \de{herunterbringen} & \luecke{1cm} \\[0.5em]
$\uparrow$ & \es{subir} \de{hochbringen} & \luecke{1cm} \\[0.5em]
$\oplus$ & \es{poner} \de{stellen, legen} & \luecke{1cm} \\
\end{tabular}}

\columnbreak

{\footnotesize
\begin{tabular}{@{}cll@{}}
$\ominus$ & \es{sacar} \de{herausnehmen} & \luecke{1cm} \\[0.5em]
$\leftarrow$ & \es{traer} \de{bringen (her)} & \luecke{1cm} \\[0.5em]
$\rightarrow$ & \es{llevar} \de{bringen (hin)} & \luecke{1cm} \\
\end{tabular}}

\end{multicols}

\vspace{2pt}

{\footnotesize \regel{Merke} \es{bajar de} / \es{sacar de} / \es{traer de} $\leftrightarrow$ \es{subir a} / \es{llevar a} / \es{poner en}}

\end{grammatikbox}

\vspace{4pt}

% ═══════════════════════════════════════════════════════════════
% KASTEN G2: OBJEKTPRONOMEN
% ═══════════════════════════════════════════════════════════════

\begin{grammatikbox}
\sanstext{\footnotesize\textbf{Kasten G2: Los pronombres de objeto directo}}

\vspace{4pt}

{\small Ergänze die Pronomen:}

\vspace{3pt}

{\footnotesize
\begin{tabular}{@{}llll@{}}
\textbf{Singular} & & \textbf{Plural} & \\
\luecke{1.5cm} \de{ihn} & $\leftarrow$ maskulin $\rightarrow$ & \luecke{1.5cm} \de{sie (m)} \\[0.4em]
\luecke{1.5cm} \de{sie (f)} & $\leftarrow$ feminin $\rightarrow$ & \luecke{1.5cm} \de{sie (f)} \\
\end{tabular}}

\vspace{4pt}

{\footnotesize \regel{Stellung} Das Pronomen steht \textbf{vor} dem konjugierten Verb: \es{La mesa} $\rightarrow$ \es{La ponemos en el salón.}}

\end{grammatikbox}

\vspace{4pt}

% ═══════════════════════════════════════════════════════════════
% KASTEN R1: REDEMITTEL
% ═══════════════════════════════════════════════════════════════

\begin{vokabelbox}
\sanstext{\footnotesize\textbf{Kasten R1: Redemittel -- La mudanza}}

\vspace{3pt}

{\footnotesize
\begin{tabular}{@{}ll@{}}
\es{¿Dónde ponemos el/la...?} & \de{Wo stellen wir den/die... hin?} \\
\es{Lo/La ponemos en...} & \de{Wir stellen ihn/sie in...} \\
\es{¿Dónde está el/la...?} & \de{Wo ist der/die...?} \\
\es{Está al lado de / detrás de / en...} & \de{Er/Sie ist neben / hinter / in...} \\
\end{tabular}}
\end{vokabelbox}

\vspace{6pt}

% ═══════════════════════════════════════════════════════════════
% AUFGABEN
% ═══════════════════════════════════════════════════════════════

\aufgabe{Aufgabe 1}{Ordne die Bilder den Vokabeln zu.} \punktebox{6}

\vspace{2pt}

{\small Verbinde jedes Symbol mit dem passenden spanischen Wort.}

\vspace{4pt}

\begin{multicols}{3}
{\footnotesize
\begin{tabular}{@{}cl@{}}
1. $\square$ Bett & a) el armario \\
2. $\square$ Lampe & b) la cama \\
3. $\square$ Kühlschrank & c) el sofá \\
4. $\square$ Sofa & d) la lámpara \\
5. $\square$ Schrank & e) el televisor \\
6. $\square$ Fernseher & f) la nevera \\
\end{tabular}

\columnbreak

\begin{tabular}{@{}cl@{}}
7. $\square$ Badezimmer & g) el salón \\
8. $\square$ Küche & h) el dormitorio \\
9. $\square$ Wohnzimmer & i) la cocina \\
10. $\square$ Schlafzimmer & j) el baño \\
11. $\square$ Flur & k) el balcón \\
12. $\square$ Balkon & l) el pasillo \\
\end{tabular}

\columnbreak

\vspace{8pt}
{\scriptsize\textit{Selbstkontrolle:}\\
Vergleiche mit Kasten V1.}

}
\end{multicols}

\vspace{8pt}

% ═══════════════════════════════════════════════════════════════

\aufgabe{Aufgabe 2}{Ergänze die Verben und Präpositionen.} \punktebox{8}

\vspace{2pt}

{\small Setze das passende Verb (konjugiert) und die richtige Präposition ein.}

\vspace{4pt}

\begin{teil}
    \item Los chicos de la mudanza \luecke{2cm} los muebles \luecke{1cm} camión.\\
    \de{Die Umzugshelfer bringen die Möbel aus dem LKW.}

    \item Después \luecke{2cm} las cajas \luecke{1cm} nuevo piso.\\
    \de{Danach bringen sie die Kisten in die neue Wohnung.}

    \item Clara \luecke{2cm} el póster \luecke{1cm} la pared.\\
    \de{Clara hängt das Poster an die Wand.}

    \item Mamá \luecke{2cm} la ropa \luecke{1cm} la caja.\\
    \de{Mama nimmt die Kleidung aus der Kiste.}
\end{teil}

\vspace{8pt}

% ═══════════════════════════════════════════════════════════════

\aufgabe{Aufgabe 3}{Ersetze das Nomen durch ein Pronomen.} \punktebox{6}

\vspace{2pt}

{\small Schreibe den Satz mit \es{lo}, \es{la}, \es{los} oder \es{las} neu.}

\vspace{4pt}

\begin{teil}
    \item \es{Ponemos la lámpara en el salón.} $\rightarrow$ \luecke{5cm}
    \item \es{Subimos los muebles al piso.} $\rightarrow$ \luecke{5cm}
    \item \es{Llevamos el sofá al dormitorio.} $\rightarrow$ \luecke{5cm}
    \item \es{Sacamos las plantas del camión.} $\rightarrow$ \luecke{5cm}
    \item \es{Bajamos la cama del piso.} $\rightarrow$ \luecke{5cm}
    \item \es{Traemos el televisor del salón.} $\rightarrow$ \luecke{5cm}
\end{teil}

\vspace{8pt}

% ═══════════════════════════════════════════════════════════════

\aufgabe{Aufgabe 4}{Beschreibt gemeinsam einen Umzug.} \punktebox{10}

\vspace{2pt}

{\small Arbeitet zu zweit. Fragt und antwortet über einen Umzug. Verwendet die Redemittel aus Kasten R1.}

\vspace{4pt}

{\footnotesize\textbf{Partner A fragt:}}

\vspace{2pt}

\begin{tabular}{@{}ll@{}}
\es{¿Dónde ponemos} & \luecke{3cm}\es{?} \\
\es{¿Y dónde está} & \luecke{3cm}\es{?} \\
\end{tabular}

\vspace{4pt}

{\footnotesize\textbf{Partner B antwortet:}}

\vspace{2pt}

\luecke{8cm}

\luecke{8cm}

\vspace{6pt}

{\footnotesize\textbf{Jetzt wechselt die Rollen! Partner B fragt:}}

\vspace{2pt}

\begin{tabular}{@{}ll@{}}
\es{¿Dónde ponemos} & \luecke{3cm}\es{?} \\
\es{¿Qué hay en} & \luecke{3cm}\es{?} \\
\end{tabular}

\vspace{4pt}

{\footnotesize\textbf{Partner A antwortet:}}

\vspace{2pt}

\luecke{8cm}

\luecke{8cm}

% ═══════════════════════════════════════════════════════════════
% FOOTER
% ═══════════════════════════════════════════════════════════════

\vfill

\noindent\rule{\textwidth}{0.5pt}

\vspace{4pt}

\hfill\sanstext{\textbf{Punkte:}} \luecke{1.2cm} / \maxpunkte

\end{document}
